% 387_Abweichungen_De_ch.tex
% Johann Pascher, 2. Okt. 2026
% Die Abweichungen im Überblick: Verhältnisse, SI-Übersetzung, Korrekturgrößen

\providecommand{\BM}{\textbf{[B]}}
\providecommand{\KM}{\textbf{[K]}}
\providecommand{\SM}{\textbf{[S]}}
\providecommand{\XM}{\textbf{[X]}}
\providecommand{\QM}{\textbf{[Q]}}

\chapter*{Die Abweichungen im Überblick}
\addcontentsline{toc}{chapter}{Die Abweichungen im Überblick}

\begin{abstract}
Dieses Dokument hält alle Restabweichungen der FFGFT an einer Stelle fest, jeweils mit dem
Messwert und seiner Unsicherheit. Es folgt der Ordnung der Theorie: zuerst die Verhältnisse der
Grundform in natürlichen Einheiten mit $\alpha=1$, dann die Übersetzung in SI. Die Verhältnisse
treffen auf Promille bis wenige Prozent; die Messunsicherheiten sind dabei fast überall um
Größenordnungen kleiner als die Abweichungen und tragen nur einen minimalen Betrag bei. Drei
Korrekturgrößen tragen die Reste: die Higgs-Korrektur $\beta_T$, die schon in den frühesten
Dokumenten von 2025 als kleine Abweichung von $\beta_T=1$ stand, die fraktale Korrektur
$K_{\text{frak}}=74/75$ der SI-Übersetzung und der Faktor~10 in $v/E_P$. Für jede wird gezeigt,
welche Varianten es gibt, wie gut sie passen und welche Zusammenhänge sich ergeben.
\end{abstract}

% ============================================================
\section*{Statusmarkierungen}
% ============================================================
\begin{center}
\begin{tabular}{@{}lp{11cm}@{}}
\textbf{[B]} & Algebraisch bewiesen. \\[2pt]
\textbf{[K]} & Numerisch bestätigt (Prüfskript, Vergleich mit Messwerten). \\[2pt]
\textbf{[S]} & Setzung, offen oder spekulativ. \\[2pt]
\textbf{[X]} & Widerlegt oder nicht tragfähig. \\[2pt]
\textbf{[Q]} & Konvention. \\
\end{tabular}
\end{center}
Messwerte: CODATA 2022, PDG 2024, Neutrinodaten wie in Dok.~340 mit den Unsicherheiten nach PDG 2024. Die Spalte
\glqq Messfehler\grqq{} nennt die relative Unsicherheit des Vergleichswerts; wo sie mit der
Abweichung vergleichbar ist, steht die Abweichung zusätzlich in Standardabweichungen.
Prüfskript: \texttt{pruef\_387\_abweichungen.py}.

% ============================================================
\section{Lesart}
% ============================================================

Die FFGFT ist in erster Linie verhältnisbasiert (Dok.~384, 386). In der Grundform, in natürlichen
Einheiten mit $\alpha=1$ und umdefinierter Ladung, gibt es nur $\xi=1/7500$ und gesetzte ganze
Zahlen; Verhältnisse prüfen die Struktur. Absolutwerte in SI brauchen einen Anker und eine
Bezugsenergie und prüfen die Umrechnung (A130). Die Abweichungen werden deshalb getrennt
aufgeführt. Eine Abweichung, die viel größer ist als der Messfehler, ist ein echter Rest der
Theorie, kein Messproblem; das gilt für fast alle Einträge.

% ============================================================
\section{Verhältnisse der Grundform}
% ============================================================

\begin{center}
\small
\begin{tabular}{@{}p{3.4cm}p{3.0cm}p{2.6cm}p{2.0cm}p{2.4cm}@{}}
\toprule
\textbf{Größe} & \textbf{FFGFT} & \textbf{Abweichung} & \textbf{Messfehler} & \textbf{Status} \\
\midrule
$m_\mu/m_e$ & $120\sqrt3=207{,}85$ & $+0{,}52\,\%$ & $2\times10^{-8}$ & \KM, Rest echt \\
$m_\tau/m_\mu$ & $\tfrac{125}{144}\xi^{-1/3}=16{,}99$ & $+1{,}03\,\%$ & $5\times10^{-5}$ & \KM, Rest echt \\
Koide $Q-\tfrac23$ & $0$ & $-0{,}017\,\xi$ & $0{,}038\,\xi$ & $0{,}4\sigma$ \\
$v/E_P$ & $\xi^4/(5\pi)$ & $-0{,}23\,\%$ & $1\times10^{-5}$ & \KM, Faktor~10 \\
$\xi$ aus $m_h/v$ & $\beta_T=0{,}977$ & $-2{,}3\,\%$ & $1{,}8\times10^{-3}$ & Prüfung \KM \\
$\Delta m^2_{\text{atm}}/\Delta m^2_{\text{sol}}$ & $120/(11/3)=32{,}7$ & $-1{,}4\,\%$ & $2{,}6\,\%$ & $0{,}5\sigma$ \\
$\sin^2\theta_{13}$ & $(25\xi)^{2/3}=0{,}0223$ & $+1{,}4\,\%$ & $3\,\%$ & $0{,}4\sigma$ \\
$M_W/M_Z$ & $\sqrt{7/9}$ & $+0{,}06\,\%$ & $1{,}7\times10^{-4}$ & $3{,}8\sigma$ \\
$m_h/M_Z$ & $11/8$ & $+0{,}15\,\%$ & $9\times10^{-4}$ & $1{,}7\sigma$ \\
$m_t/M_Z$ & $(11/8)^2$ & $-0{,}10\,\%$ & $1{,}7\times10^{-3}$ & $0{,}6\sigma$ \\
$\lambda_{\text{CKM}}$ & $\xi^{1/6}=0{,}2260$ & $+0{,}46\,\%$ & $3\times10^{-3}$ & $1{,}5\sigma$ \\
\bottomrule
\end{tabular}
\end{center}

Die Leptonverhältnisse sind die schärfsten Prüfungen der Grundform. Ihre Messfehler liegen bei
$10^{-8}$ und $10^{-5}$, die Abweichungen bei $0{,}5$ und $1\,\%$; die Reste sind also vollständig
Reste der Leiter, die Messung trägt nichts bei \KM. Ihre Ursache ist offen (R113). Bei den
Neutrinos, bei $\theta_{13}$, $m_t$ und $\lambda_{\text{CKM}}$ liegen die Abweichungen innerhalb
von 1,5 Standardabweichungen; dort ist die Messung der begrenzende Faktor. $M_W$ aus
$\sin^2\theta_W=2/9$ ist mit $3{,}8\sigma$ die einzige ausgewiesene Spannung in der Grundform.
Koide trifft innerhalb der Messgenauigkeit von $m_\tau$.

% ============================================================
\section{Die Übersetzung in SI}
% ============================================================

\begin{center}
\small
\begin{tabular}{@{}p{3.6cm}p{3.0cm}p{2.4cm}p{2.0cm}p{2.4cm}@{}}
\toprule
\textbf{Größe} & \textbf{FFGFT} & \textbf{Abweichung} & \textbf{Messfehler} & \textbf{Status} \\
\midrule
$\alpha^{-1}$ mit $74/75$ & $137{,}059$ & $+1{,}7\times10^{-4}$ & $1{,}5\times10^{-10}$ & \KM, Anker \\
$\alpha^{-1}$ Galois & $3700/27=137{,}037$ & $+7{,}6\times10^{-6}$ & $1{,}5\times10^{-10}$ & \KM \\
$m_em_\mu$ gegen $54$\,MeV$^2$ & $54$ & $-1{,}6\times10^{-4}$ & $2\times10^{-8}$ & \KM \\
$m_e$ (Kette, $E_P$ gemessen) & $\tfrac{4}{15\pi}\xi^{11/2}E_P$ & $-1{,}32\,\%$ & $1\times10^{-5}$ & \KM \\
$m_\mu$ (Kette) & $\tfrac{16}{25\pi}\xi^{5}E_P$ & $-0{,}80\,\%$ & $1\times10^{-5}$ & \KM \\
$m_\tau$ (Kette) & $\tfrac{5}{9\pi}\xi^{14/3}E_P$ & $+0{,}22\,\%$ & $5\times10^{-5}$ & \KM \\
$v$ aus Leiter, nackt & $248{,}93$\,GeV & $+1{,}10\,\%$ & $10^{-6}$ & \KM \\
$v$ aus Leiter mit $K$ & $245{,}61$\,GeV & $-0{,}25\,\%$ & $10^{-6}$ & \KM \\
$\Delta m^2_{\text{atm}}$ & $2{,}476\times10^{-3}$\,eV$^2$ & $-1{,}0\,\%$ & $1{,}1\,\%$ & $0{,}9\sigma$ \\
$\Delta m^2_{\text{sol}}$ & $7{,}565\times10^{-5}$\,eV$^2$ & $+0{,}46\,\%$ & $2{,}4\,\%$ & $0{,}2\sigma$ \\
\bottomrule
\end{tabular}
\end{center}

Auch hier sind die Messfehler bis auf die Neutrinos um Größenordnungen kleiner als die
Abweichungen. Zwei Reste fallen auf. Der Rest von $\alpha$, $1{,}7\times10^{-4}$, ist genau der
Abstand der gemessenen Massen vom Galois-Wert: $m_em_\mu=53{,}991$\,MeV$^2$ liegt
$1{,}6\times10^{-4}$ unter $54=|\mathrm{GF}(3)^*|\cdot|\mathrm{GF}(27)|$, und mit $54$ trifft
$\alpha^{-1}=3700/27$ auf $7{,}6$\,ppm \KM. Der Rest der Elektronmasse in der Kette, $-1{,}32\,\%$,
ist bis auf $1{,}6\times10^{-4}$ gleich $K_{\text{frak}}-1=-1{,}33\,\%$ \KM; für Myon und Tau gilt
das nicht.

% ============================================================
\section{Die Korrekturgrößen}
% ============================================================

\subsection{\texorpdfstring{$\beta_T$}{beta\_T}: die Higgs-Korrektur aus den frühesten Dokumenten}

Die frühesten Dokumente vom März und April 2025 setzen zwei Größen auf~1: $\alpha=1$ und
$\beta_T=1$, mit der Begründung, dass beide in natürlichen Einheiten Konventionen sind und ihre
SI-Werte, $\alpha\approx1/137$ und $\beta\approx0{,}008$, erst durch die Übersetzung entstehen
(\glqq Die Konsistenz von $\alpha=1$ und $\beta=1$\grqq, Versionsgeschichte). $\beta_T=1$ mit
$r_0=\xi\,\ell_P$ ergab die Higgs-Formel $\xi=\lambda_h^2v^2/(16\pi^3m_h^2)$; das Dokument nannte
dabei $r_0\approx\ell_P/7519$ statt $\ell_P/7500$, also schon damals eine kleine Abweichung von
$\beta_T=1$ (Dok.~385). Heute gerechnet ist
\begin{equation}
  \beta_T=\frac{\xi_{\text{EFT}}}{\xi}=\frac{m_h^2}{64\pi^3v^2\,\xi}=0{,}977\quad\KM ,
\end{equation}
mit einem Messfehler von $0{,}18\,\%$ aus $m_h$. Die Abweichung von $\beta_T=1$ ist damit ein
echter Rest, 13-mal größer als die Messunsicherheit. Sie hängt davon ab, welches $v$ eingesetzt
wird:
\begin{center}
\small
\begin{tabular}{@{}lcc@{}}
\toprule
\textbf{Eingabe für $v$} & $v$ & $\beta_T$ \\
\midrule
gemessen (aus $G_F$) & $246{,}22$\,GeV & $0{,}977$ \\
Kette $\xi^4E_P/(5\pi)$ & $245{,}65$\,GeV & $0{,}982$ \\
Leiter mit $m_e$ und $K_{\text{frak}}$ & $245{,}61$\,GeV & $0{,}982$ \\
Leiter mit $m_e$, nackt & $248{,}93$\,GeV & $0{,}956$ \\
\bottomrule
\end{tabular}
\end{center}
Mit den FFGFT-eigenen Werten für $v$ rückt $\beta_T$ näher an~1, bleibt aber bei
$-1{,}8\,\%$; mit dem nackten $v$ entfernt es sich. Die Higgs-Prüfung bleibt eine Prüfung von
$\xi$ auf wenige Prozent \KM. Eine Herleitung des Rests gibt es nicht \SM. Numerisch liegt
$\beta_T=0{,}9772$ nahe bei $K_{\text{frak}}^{\sqrt3}=0{,}9770$; bei einem Messfehler von
$0{,}18\,\%$ ist das keine belastbare Aussage.

\subsection{\texorpdfstring{$K_{\text{frak}}$}{K\_frak}: die Korrektur der SI-Übersetzung}

In der Grundform gibt es keinen Korrekturfaktor; die Brücke lautet $\xi E_0^2=1$ (Dok.~386). In
der SI-Übersetzung verlangt $\alpha$ den Faktor $K_\alpha=\xi m_em_\mu\alpha^{-1}=0{,}98650$. Die
Varianten im Vergleich:
\begin{center}
\small
\begin{tabular}{@{}lccccc@{}}
\toprule
\textbf{Variante} & $K$ & $\alpha^{-1}$ & \textbf{Galois} & $v$ (Leiter) & \textbf{Umläufe} \\
\midrule
$1-100\xi=74/75$ & $0{,}986667$ & $+1{,}7\times10^{-4}$ & $+7{,}6$\,ppm & $-0{,}25\,\%$ & $75{,}0$ \\
$(1-\xi)^{100}\approx e^{-100\xi}$ & $0{,}986754$ & $+2{,}6\times10^{-4}$ & $+96$\,ppm & $-0{,}24\,\%$ & $75{,}5$ \\
$1/(1+100\xi)$ & $0{,}986842$ & $+3{,}5\times10^{-4}$ & $+190$\,ppm & $-0{,}23\,\%$ & $76{,}0$ \\
$1-\psi'(75)$ (Dok.~381) & $0{,}986577$ & $+7{,}8\times10^{-5}$ & $-83$\,ppm & $-0{,}26\,\%$ & $74{,}5$ \\
ohne Korrektur & $1$ & $+1{,}4\,\%$ & $+1{,}4\,\%$ & $+1{,}10\,\%$ & -- \\
\bottomrule
\end{tabular}
\end{center}
Nur $74/75$ schließt die Windung nach einer ganzen Zahl von Umläufen und trifft die Galois-Form
auf $7{,}6$\,ppm \KM. $1-\psi'(75)$ halbiert den $\alpha$-Rest, verliert aber die ganze Umlaufzahl
und verschlechtert die Galois-Form. Die multiplikative und die inverse Form passen überall
schlechter als die additive; zur offenen Form (R139) spricht der $\alpha$-Sektor also für
$1-100\xi$. Auf die Leptonverhältnisse wirkt keine Variante, auf $v$ und $G$ unterscheiden sie
sich nur um $10^{-4}$.

\subsection{Der Faktor~10 in \texorpdfstring{$v/E_P$}{v/E\_P}}

$v/E_P=\xi^4/(5\pi)$ mit $5\pi=(\pi/2)\cdot10$ trifft auf $-0{,}23\,\%$ bei einem Messfehler von
$10^{-5}$. Der Faktor~10 ist an den Messwert angepasst: Im Arbeitsstand vom Februar 2026 ergab
$m_P/(f^4\cdot\pi/2)$ den Wert $2467$\,GeV statt $246$\,GeV, und die Lücke von $10{,}02$ wurde mit
$1/10$ geschlossen. Zwei Lesarten wurden geprüft:
\begin{itemize}[leftmargin=*,itemsep=1pt]
\item Rationalisierte Gravitation ($8\pi G=1$) mit der reduzierten Planck-Masse ergäbe
  $2\sqrt{8\pi}=10{,}03$ statt~10. Dok.~374 und 376 zeigen aber, dass die passende
  Referenzenergie die nicht reduzierte Planck-Energie ist; die reduzierte Lesart wäre eine andere
  physikalische Annahme. Diese Lesart trägt nicht \XM.
\item Gemessenes gegen nacktes $v$: Setzt man statt des gemessenen $v=246{,}22$\,GeV das nackte
  $v=248{,}93$\,GeV der Leiter ein, verlangt die Formel den Nenner $15{,}501$. Das ist
  $5\pi\,K_{\text{frak}}=15{,}4985$ oder $\pi^3/2=15{,}5031$, jeweils auf $1{,}5\times10^{-4}$. Der
  Faktor wäre dann $10\,K_{\text{frak}}\approx\pi^2$ statt~10. Zwischen beiden lässt sich auf
  diesem Niveau nicht unterscheiden \SM.
\end{itemize}
Die zweite Lesart passt zur Trennung von Grundform und SI: Der Faktor~10 wäre dann eine
Vermischung des nackten Werts der Leiter mit dem SI-Vergleichswert von $v$. Eine Herleitung ist
das nicht.

\textit{(Vermerk vom 2.~Okt.~2026 zur Bezeichnung: $f$ ist hier die Grundwindungszahl $f = 1/\xi = 7500$, keine Frequenz. In natürlichen Einheiten ist $f$ nur eine andere Lesart der Dualität $T\cdot m = 1$: $f = t_P/t_0 = \Lambda_{\text{T0}}/E_P$ mit $t_0 = \xi\,t_P$ und $\Lambda_{\text{T0}} = E_P/\xi$, und aus $\xi E_0^2 = 1$ folgt $f = E_0^2$. Als dimensionsloses Verhältnis lässt sich $f$ nicht begründet auf~1 setzen (Dok.~A135). Vgl.\ Dok.~190, R147.)}

% ============================================================
\section{Zusammenhänge}
% ============================================================

Die Reste stehen nicht unverbunden nebeneinander:
\begin{enumerate}[leftmargin=*,itemsep=2pt]
\item \textbf{$\alpha$ und Galois.} Der $\alpha$-Rest von $1{,}7\times10^{-4}$ ist der Abstand
  $m_em_\mu$ gegen $54$\,MeV$^2$. Mit dem Galois-Wert bleibt $K_{\text{frak}}=74/75$ exakt, und der
  Rest liegt auf der Seite der Massen \KM.
\item \textbf{Elektron und $K_{\text{frak}}$.} Der Elektronrest der Kette ist $K_{\text{frak}}-1$
  bis auf $1{,}6\times10^{-4}$. Die Kette enthält aber den Faktor~10; solange er offen ist, ist
  diese Übereinstimmung nicht belastbar \SM.
\item \textbf{Faktor~10 und $K_{\text{frak}}$.} Mit dem nackten $v$ wird der Faktor~10 zu
  $10\,K_{\text{frak}}$; dieselbe Korrektur, die $E_0$ und $\alpha$ trägt, stünde dann auch in
  $v/E_P$ \SM.
\item \textbf{$\beta_T$.} Die Higgs-Korrektur ist mit $-2{,}3\,\%$ die größte und älteste. Mit den
  FFGFT-eigenen Werten für $v$ sinkt sie auf $-1{,}8\,\%$; einen Zusammenhang mit
  $K_{\text{frak}}$ zeigt sie nicht in belastbarer Form \SM.
\end{enumerate}

% ============================================================
\section{Messgenauigkeit}
% ============================================================

\begin{center}
\small
\begin{tabular}{@{}lll@{}}
\toprule
\textbf{Eingang} & \textbf{relative Unsicherheit} & \textbf{Wirkung} \\
\midrule
$\alpha$, $m_e$ & $1{,}5\times10^{-10}$, $3\times10^{-10}$ & vernachlässigbar \\
$m_\mu$, $M_Z$ & $2\times10^{-8}$, $2\times10^{-5}$ & vernachlässigbar \\
$G$ (also $E_P$) & $2\times10^{-5}$ ($1\times10^{-5}$) & vernachlässigbar \\
$m_\tau$ & $5\times10^{-5}$ & vernachlässigbar außer Koide \\
$m_h$ & $9\times10^{-4}$ & $\beta_T$ auf $0{,}18\,\%$ \\
$m_t$, $M_W$ & $1{,}7\times10^{-3}$, $1{,}7\times10^{-4}$ & begrenzend \\
Neutrinodaten, $\theta_{13}$, $\lambda_{\text{CKM}}$ & $1$ bis $3\,\%$ & begrenzend \\
\bottomrule
\end{tabular}
\end{center}
Für die Leptonmassen, $\alpha$, $v/E_P$ und die SI-Kette tragen die Messfehler nur einen minimalen
Betrag bei; die Reste dort sind Reste der Theorie. Für $m_h$, $m_t$, $M_W$, die Neutrinos und die
Mischungswinkel ist die Messung der begrenzende Faktor, und die Abweichungen stehen dort in
Standardabweichungen.

% ============================================================
\section{Fazit}
% ============================================================

Die Verhältnisse der Grundform treffen auf Promille bis wenige Prozent, und wo die Messung genau
ist, sind die Reste echte Reste der Theorie. Drei Korrekturgrößen tragen sie. $\beta_T$ ist die
älteste: Schon die frühesten Dokumente setzten $\beta_T=1$ neben $\alpha=1$ und sahen eine kleine
Abweichung; heute ist es $0{,}977$, mit FFGFT-eigenem $v$ $0{,}982$. $K_{\text{frak}}=74/75$ ist
die Korrektur der SI-Übersetzung; unter den Varianten ist sie die einzige mit ganzer Umlaufzahl
und trifft die Galois-Form auf $7{,}6$\,ppm. Der Faktor~10 in $v/E_P$ ist angepasst; mit dem
nackten $v$ der Leiter wird er zu $10\,K_{\text{frak}}\approx\pi^2$. Offen sind die Ursache der
Leptonreste, der Rest von $\beta_T$ und die Herkunft des Faktors in $v/E_P$.

\vspace{1em}
\noindent\textit{Prüfskript:} \texttt{2/python/Dok387\_Skripte/pruef\_387\_abweichungen.py}
--- 27/27.
